\documentclass[10pt,a4paper]{amsart}
\usepackage[margin=19mm]{geometry}
\usepackage{amsmath,amssymb,mathtools}
\usepackage[scr=rsfs,cal=euler]{mathalfa}
\usepackage{xfrac}
\usepackage[unicode,hidelinks,bookmarks=false]{hyperref}
\newcommand{\ie}{i.e.\ }
\newcommand{\cf}{cf.\ }
\newcommand{\resp}{respectively\ }
\newcommand{\Subsection}{\subsection}
\providecommand{\nobreakdash}{\nobreak\hbox{-}}
\setlength{\emergencystretch}{2em}
\multlinegap0pt
\usepackage{bm}
\usepackage{bbm}
\usepackage{dsfont}
\usepackage{xspace}
\usepackage{cancel}
\usepackage{mathtools}
\usepackage[shortlabels]{enumitem}
\usepackage{amsthm}
\makeatletter
\@ifundefined{theorem}{\newtheorem{theorem}{Theorem}}{}
\@ifundefined{definition}{\newtheorem{definition}[theorem]{Definition}}{}
\makeatother
\title[The Subtractive Divisor Orbit: Unconditional Bounds, Parity ]{The Subtractive Divisor Orbit: Unconditional Bounds, Parity Constraints, and a Conditional Framework}
\author{Marco Mantovanelli}
\date{Source: arXiv:2604.25446v2 (CC BY 4.0)}
\begin{document}
\maketitle

\noindent\emph{Source and licence.} The Subtractive Divisor Orbit: Unconditional Bounds, Parity Constraints, and a Conditional Framework. Version arXiv:2604.25446v2 (CC BY 4.0), distributed under the Creative Commons Attribution 4.0 International licence, as recorded in the arXiv licence metadata for this exact version and shown on the abstract page https://arxiv.org/abs/2604.25446v2. Full rights notice: licenses/arxiv-2604.25446-CC-BY-4.0.txt.\par\smallskip

\noindent\textit{Editorial scope.} Counted unit: the Theorem whose source label is \texttt{thm:variance-ladder} in the article (source0.tex lines 580--601, source label \texttt{thm:variance-ladder}), delivered together with the article's own complete original proof of it (source0.tex lines 603--650, one \texttt{proof} environment of 48 lines). No mathematics is written, repaired or re-derived: statement and proof are the article's own text line for line. Required context retained uncounted: def:dynamic-ladder (lines 557--573). Every definition, display and label of those lines is retained unchanged. Omitted from this package and not redistributed: 1--556, 574--579, 602--602, 651--988. Nothing inside a retained excerpt is omitted or altered: every retained body is delivered exactly as the article's lines are written, so no omission receipt of any kind accompanies this package. Citations: the retained excerpts carry no citation command at all, so no bibliography entry of the article is transcribed and no reference is redistributed. Reference adaptation: none was needed and none was made; every reference inside the delivered unit resolves inside it, so no unresolved reference or citation is produced. Printed slips kept literal: no printed word, symbol, hypothesis or number of the counted statement or of its proof was altered. Layout and class: the article's own class is \texttt{article} with packages including fontenc, lmodern, amsmath, amssymb, amsthm, mathtools, geometry, microtype, enumitem, booktabs, biblatex, hyperref; the wrapper is the lane's pinned \texttt{amsart} editorial wrapper at 10pt on A4, which restates the theorem-like environments the retained text needs and adds the article's own macro definitions that the retained text needs (none), with extra wrapper packages (bm, bbm, dsfont, xspace, cancel, mathtools, enumitem). The delivered numbering is the unit's own. Assets: no retained span contains a figure environment, a graphics-inclusion command, a PDF-page inclusion, a table or a TikZ picture; no image file is copied into this package, the package declares no asset dependency and no image file is redistributed with it. Licence: version arXiv:2604.25446v2 of this article is distributed under the Creative Commons Attribution 4.0 International licence, as recorded in the arXiv licence metadata for this exact version and shown on the abstract page https://arxiv.org/abs/2604.25446v2. A full rights notice is delivered at licenses/arxiv-2604.25446-CC-BY-4.0.txt. Characters: every retained excerpt is pure ASCII, so the delivered engine, XeLaTeX with its default Latin Modern text font, reports no missing character.
\begin{definition}[Dynamic near-ladder]\label{def:dynamic-ladder}
Let $0<\eta<1$.  An ordered subsequence
\[
 m_1>m_2>\cdots>m_r
\]
of states from $J_x(N)$ is an $\eta$-\emph{dynamic near-ladder} with
parameter $T$ if
\begin{enumerate}[label=\textup{(\roman*)}]
 \item $|\tau(m_i)-T|\leq\eta T$ for every $i$;
 \item $m_i-m_{i+1}=\tau(m_i)+s_i$ with $s_i\geq0$;
 \item $\sum_{i<r}s_i\leq\eta E$.
\end{enumerate}
It is \emph{step-mass-saturating} if, in addition,
\[
 \sum_{i=1}^r\tau(m_i)\geq(1-\eta)E.
\]
\end{definition}

\begin{theorem}[Small variance produces a dynamic ladder]
\label{thm:variance-ladder}
Assume $0<\varepsilon\leq1/16$ and
\begin{equation}\label{eq:variance-assumption}
 \frac1V\sum_{j\in J}(\tau(n_j)-T)^2
 \leq\varepsilon^2T^2.
\end{equation}
Then the states satisfying
\[
 |\tau(n_j)-T|\leq\sqrt\varepsilon\,T
\]
form a $3\sqrt\varepsilon$-dynamic near-ladder which is
step-mass-saturating.  Moreover, for all but
$O(\sqrt\varepsilon\,r)$ adjacent ladder pairs,
\begin{equation}\label{eq:ladder-gaps}
 m_i-m_{i+1}=T+O(\sqrt\varepsilon\,T).
\end{equation}
In particular,
\[
 r=(1+O(\varepsilon))V\asymp\frac ET.
\]
\end{theorem}

\begin{proof}
Let
\[
 \mathcal E=\{j\in J:|\tau(n_j)-T|>\sqrt\varepsilon\,T\}.
\]
Chebyshev's inequality and \eqref{eq:variance-assumption} give
\begin{equation}\label{eq:exception-count}
 |\mathcal E|\leq\varepsilon V.
\end{equation}
The step mass of this exceptional set is bounded by
\begin{align*}
 \sum_{j\in\mathcal E}\tau(n_j)
 &\leq T|\mathcal E|
 +|\mathcal E|^{1/2}
   \left(\sum_{j\in J}(\tau(n_j)-T)^2\right)^{1/2}\\
 &\leq (\varepsilon+\varepsilon^{3/2})TV
 \leq2\varepsilon E.
\end{align*}
Thus the complementary regular set carries at least $(1-2\varepsilon)E$
step mass.  List its states in orbit order as $m_1>\cdots>m_r$.
Equation \eqref{eq:exception-count} gives
$r=(1+O(\varepsilon))V$.

Between consecutive regular states, the orbit first makes the step
$\tau(m_i)$ and then possibly passes through exceptional states.  Hence
\[
 m_i-m_{i+1}=\tau(m_i)+s_i,
\]
where $s_i\geq0$ is the intervening exceptional step mass.  Summing in
$i$ gives
\[
 \sum_{i<r}s_i\leq2\varepsilon E.
\]
Because $\varepsilon\leq1/16$, all requirements of
Definition~\ref{def:dynamic-ladder} hold with
$\eta=3\sqrt\varepsilon$.

Finally, the number of indices with
$s_i>\sqrt\varepsilon T$ is at most
\[
 \frac{2\varepsilon E}{\sqrt\varepsilon T}
 =2\sqrt\varepsilon V
 =O(\sqrt\varepsilon r).
\]
For the remaining pairs, combine this with
$\tau(m_i)=T+O(\sqrt\varepsilon T)$ to obtain
\eqref{eq:ladder-gaps}.
\end{proof}

\end{document}
